\clearpage
\seccionInforme{7. Anexos}
\phantomsection\addcontentsline{toc}{section}{7. Anexos}

El código fuente completo de la aplicación está disponible en el repositorio \url{https://github.com/OwerLopez/matrilab-epis-react-native}. En este anexo se reúnen capturas complementarias.

\vspace{0.2cm}
\phantomsection\addcontentsline{toc}{subsection}{Anexo A. Tema oscuro}
\colabtitle{Anexo A. Tema oscuro}

Al elegir el tema oscuro desde el perfil, toda la aplicación cambia de paleta: fondos, tarjetas, textos y los colores del semáforo usan sus variantes oscuras manteniendo el contraste (Figura \ref{fig:oscuro}). En el detalle (c) se ve otra regla de negocio: el estudiante ya tiene un grupo de Plataformas Emergentes y no puede tomar un segundo grupo del mismo curso.

\begin{figure}[!htbp]
\centering
\safeimage{1.0\textwidth}{img/figuras/fig_oscuro.png}{Tema oscuro: (a) panel del estudiante, (b) catálogo, (c) detalle con la regla de un grupo por curso, (d) horario}
\caption{Tema oscuro: (a) panel del estudiante, (b) catálogo, (c) detalle con la regla de un grupo por curso, (d) horario}
\label{fig:oscuro}
\end{figure}

\vspace{0.2cm}
\phantomsection\addcontentsline{toc}{subsection}{Anexo B. Desarrollo con Expo Go}
\colabtitle{Anexo B. Desarrollo con Expo Go}

El menú de desarrollo de Expo Go (Figura \ref{fig:entorno}, panel b) permite recargar la aplicación, activar el monitor de rendimiento, inspeccionar elementos de la interfaz y abrir las herramientas de depuración. Se usó sobre todo el inspector de elementos para revisar los márgenes y los contenedores de cada pantalla.

